% !Mode:: "TeX:UTF-8"

\chapter{总结与展望}[Conclusion and Prospect]

\section{结果讨论与总体总结}[Discussion and overall summary]

综合第 5 章实验结果可以看出，本文方法在六类代表性动作上总体实现了较为稳定的整身动作跟踪，其中周期性动作和静态平衡动作表现最好，而高动态跳跃、快速变向以及带有明显腾空落地过程的动作仍然更具挑战。误差主要集中在冲击响应、重心快速切换以及上肢与躯干的高频协同控制阶段。尽管如此，多模态参考约束、PPO 优化框架以及课程化稳定训练策略仍显著提升了整身动作的可执行性和稳定性，使机器人在复杂动作和扰动条件下保持了较高成功率与较好的泛化表现。

围绕“基于强化学习的人形机器人运动技能学习与泛化研究”这一课题，本文构建了从人体专家数据处理、动作重定向，到强化学习动作跟踪训练与实验验证的完整技术流程。针对人体运动数据无法直接用于机器人控制的问题，本文建立了 SMPL 人体模型与目标人形机器人之间的关键点拓扑映射关系，并通过体型拟合、逐帧动作重定向、关节限位约束和时序高斯滤波，将人体运动捕捉序列转换为机器人可执行的关节空间参考轨迹。在运动技能学习方面，本文基于 HumanoidVerse 和 Isaac Gym/MuJoCo 仿真环境构建整身动作跟踪任务，以重定向后的专家轨迹为参考状态，通过相位编码、多模态误差度量和高斯核奖励函数引导策略学习，并采用 PPO 算法优化控制策略。针对高动态动作训练中容易出现的局部最优和策略保守化问题，本文进一步引入参考状态初始化、动作偏离终止课程和安全正则项，以提升策略收敛效率和物理可行性。整体来看，本文提出的数据处理与强化学习训练流程能够较稳定地完成多类代表性动作的跟踪任务，并为后续更复杂的人形机器人运动技能学习研究奠定了基础。

\section{研究局限性}[Limitations]

尽管本文已经完成了从专家数据预处理到动作跟踪训练与测试的基本闭环，但当前研究仍存在若干局限。首先，本文实验主要集中在仿真环境中，尚未开展真实机器人平台验证，因此策略在真实执行器、传感器误差和复杂环境接触条件下的表现仍需进一步确认。其次，目前测试动作数量和场景范围仍然有限，虽然已经覆盖了静态平衡、周期行走和高动态动作，但对更长时序、更强交互性和更复杂任务语义条件下的动作学习研究还不充分。再次，当前奖励设计和课程学习参数仍依赖人工经验进行设定，不同动作之间的统一权重配置和自动调参机制仍有进一步优化空间。

\section{未来工作展望}[Future Work]

未来工作可以从以下几个方向继续展开。第一，进一步扩展动作类别和测试场景，构建覆盖更多全身运动技能的参考动作库，并系统评估策略在复杂地形、连续任务和长时序动作中的表现。第二，开展跨仿真平台以及真机验证研究，分析不同动力学模型、执行器延迟和接触参数差异对策略迁移性能的影响，为后续 sim-to-sim 和 sim-to-real 打下基础。第三，继续优化奖励函数设计、课程学习策略和观测建模方式，提升策略在高动态动作下的稳定性与样本效率。第四，尝试将语言指令、视觉感知或多模态高层任务条件引入现有底层运动控制框架中，推动人形机器人从单一动作模仿逐步过渡到面向开放场景的具身智能控制。
